\documentclass{standalone}
\usepackage{tikz}
\usepackage{animate}
\usepackage{graphicx}
\usetikzlibrary{fit, positioning, arrows.meta, shapes, calc}

\begin{document}
\begin{animateinline}[controls, loop]{1}
\multiframe{12}{iFrame=1+1}
{
\begin{tikzpicture}[
    >=Stealth,
    image_node/.style={
        inner sep=0pt, align=center
    },
    red_box_style/.style={
        draw=red!60!black!50, fill=red!10, thick, rounded corners=15pt
    },
    blue_box_style/.style={
        draw=blue!60!black!50, fill=blue!10, thick, rounded corners=15pt
    },
    dashed_box/.style={
        draw=black!60, thick, dashed, rounded corners=10pt, fill=white
    },
    thick_red_arrow/.style={
        draw=red!60!black!70, line width=3mm, -{Triangle[length=4mm, width=6mm]}
    },
    blue_arrow/.style={
        draw=blue!60!black, line width=1.5mm, -{Triangle[length=4mm, width=5mm]}
    }
]

    % Explicit Spacious Bounding Box
    \path [use as bounding box] (-2, -4) rectangle (28, 5);

    % ==========================================
    % 1. PARENT BLOCKS (Drawn first to stay in background)
    % ==========================================
    
    % Anchor coordinates for the Red Box (Blur kernel estimation)
    \coordinate (red_top_left) at (2.5, 3.5);
    \coordinate (red_bottom_right) at (19.0, -3.5);
    \node[fit={(red_top_left) (red_bottom_right)}, red_box_style] (red_box) {};
    \node[above=0.3cm of red_box, font=\Large\sffamily] {Blur kernel estimation};

    % Anchor coordinates for the Blue Box (Non-blind deconvolution)
    \coordinate (blue_top_left) at (20.0, 3.5);
    \coordinate (blue_bottom_right) at (27.0, -3.5);
    \node[fit={(blue_top_left) (blue_bottom_right)}, blue_box_style] (blue_box) {};
    \node[above=0.3cm of blue_box, font=\Large\sffamily] {Non-blind deconvolution};


    % ==========================================
    % 2. EXTERNAL INPUT
    % ==========================================
    
    \node[image_node] (img_input) at (0, 0) {\includegraphics[width=2.5cm, height=2.5cm, keepaspectratio]{crops/2015_07298601-Figure4-1subFig-1/input.png}};
    \node[above=0.2cm of img_input, font=\sffamily] {input $\mathbf{y}$};


    % ==========================================
    % 3. RED BOX CONTENTS
    % ==========================================
    
    % Edge Mask
    \node[image_node] (img_mask) at (4.5, 0) {\includegraphics[width=2.5cm, height=2.5cm, keepaspectratio]{crops/2015_07298601-Figure4-1subFig-1/edge_mask.png}};
    \node[above=0.2cm of img_mask, font=\sffamily] {edge mask $M$};
    % Small red highlight square on edge mask
    \draw[red, very thick] ($(img_mask.center) + (0.6, 0.2)$) rectangle +(0.3, 0.3);

    % Optimize Patch Centers (Dashed Box)
    \coordinate (opt_patch_tl) at (7.2, 2.5);
    \coordinate (opt_patch_br) at (11.8, -2.8);
    \node[fit={(opt_patch_tl) (opt_patch_br)}, dashed_box] (opt_patch_box) {};
    \node[fill=white, inner sep=2pt, font=\sffamily] at (opt_patch_box.north) {optimize patch centers};

    % Contents of Dashed Box
    \node[image_node] (scatter) at (9.5, 1.2) {\includegraphics[width=4cm, height=1.8cm, keepaspectratio]{crops/2015_07298601-Figure4-1subFig-1/upper_patch.png}};
    
    \node[image_node] (patch_p) at (8.2, -1.4) {\includegraphics[width=1.6cm, height=1.6cm, keepaspectratio]{crops/2015_07298601-Figure4-1subFig-1/p1.png}};
    \node[above=0.1cm of patch_p, font=\sffamily] {$\mathbf{P}_i \hat{\mathbf{x}}$};
    
    \node[image_node] (patch_z) at (10.8, -1.4) {\includegraphics[width=1.6cm, height=1.6cm, keepaspectratio]{crops/2015_07298601-Figure4-1subFig-1/p2.png}};
    \node[above=0.1cm of patch_z, font=\sffamily] {$\mathbf{z}_i$};
    
    \draw[blue_arrow] (patch_p) -- (patch_z);

    % Intermediate Latent Image
    \node[image_node] (latent) at (14.5, 0) {\includegraphics[width=2.5cm, height=2.5cm, keepaspectratio]{crops/2015_07298601-Figure4-1subFig-1/latent_image.png}};
    \node[above=0.2cm of latent, align=center, font=\sffamily] {intermediate\\latent image $\hat{\mathbf{x}}$};
    % Small red highlight square on latent image
    \draw[red, very thick] ($(latent.center) + (0.6, 0.2)$) rectangle +(0.3, 0.3);

    % Estimate Blur Kernel
    \node[image_node] (kernel1) at (17.5, 0) {\includegraphics[width=1.8cm, height=1.8cm, keepaspectratio]{crops/2015_07298601-Figure4-1subFig-1/kernel.png}};
    \node[above=0.2cm of kernel1, align=center, font=\sffamily] {estimate\\blur kernel\\$\mathbf{k}$};


    % ==========================================
    % 4. BLUE BOX CONTENTS
    % ==========================================
    
    % Final Blur Kernel
    \node[image_node] (kernel2) at (22.0, 0) {\includegraphics[width=1.8cm, height=1.8cm, keepaspectratio]{crops/2015_07298601-Figure4-1subFig-1/final_kernel.png}};
    \node[above=0.2cm of kernel2, align=center, font=\sffamily] {final\\blur kernel};

    % Output Image
    \node[image_node] (output) at (25.0, 0) {\includegraphics[width=2.5cm, height=2.5cm, keepaspectratio]{crops/2015_07298601-Figure4-1subFig-1/output.png}};
    \node[above=0.2cm of output, font=\sffamily] {output};


    % ==========================================
    % 5. CONNECTIONS / ARROWS
    % ==========================================
    
    % Thick Red Block Arrows
    \draw[thick_red_arrow] (img_input.east) -- (red_box.west);
    \draw[thick_red_arrow] (red_box.east) -- (blue_box.west);

    % Blue Main Pipeline Arrows
    \draw[blue_arrow] (img_mask.east) -- (opt_patch_box.west);
    \draw[blue_arrow] (opt_patch_box.east) -- (latent.west);
    \draw[blue_arrow] (latent.east) -- (kernel1.west);
    \draw[blue_arrow] (kernel2.east) -- (output.west);

    % Blue Feedback Loop Arrow
    \draw[blue_arrow] (kernel1.south) -- (17.5, -3.1) -- (9.5, -3.1) -- (opt_patch_box.south);

    % ==========================================
    % 6. CONDITIONAL HOPPING BOUNDING BOXES
    % ==========================================
    % placeholder for frame 1: input y
    \ifnum\iFrame=1
        \node[draw=red, very thick, dashed, inner sep=5pt, fit=(img_input)] {};
    \fi
    
    % placeholder for frame 2: red_box
    \ifnum\iFrame=2
        \node[draw=red, very thick, dashed, inner sep=5pt, fit=(red_box)] {};
    \fi
    
    % placeholder for frame 3: img_mask
    \ifnum\iFrame=3
        \node[draw=red, very thick, dashed, inner sep=5pt, fit=(img_mask)] {};
    \fi
    
    % placeholder for frame 4: opt_patch_box
    \ifnum\iFrame=4
        \node[draw=red, very thick, dashed, inner sep=5pt, fit=(opt_patch_box)] {};
    \fi
    
    % placeholder for frame 5: scatter
    \ifnum\iFrame=5
        \node[draw=red, very thick, dashed, inner sep=5pt, fit=(scatter)] {};
    \fi
    
    % placeholder for frame 6: patch_p
    \ifnum\iFrame=6
        \node[draw=red, very thick, dashed, inner sep=5pt, fit=(patch_p)] {};
    \fi
    
    % placeholder for frame 7: patch_z
    \ifnum\iFrame=7
        \node[draw=red, very thick, dashed, inner sep=5pt, fit=(patch_z)] {};
    \fi
    
    % placeholder for frame 8: latent
    \ifnum\iFrame=8
        \node[draw=red, very thick, dashed, inner sep=5pt, fit=(latent)] {};
    \fi
    
    % placeholder for frame 9: kernel1
    \ifnum\iFrame=9
        \node[draw=red, very thick, dashed, inner sep=5pt, fit=(kernel1)] {};
    \fi
    
    % placeholder for frame 10: blue_box
    \ifnum\iFrame=10
        \node[draw=red, very thick, dashed, inner sep=5pt, fit=(blue_box)] {};
    \fi
    
    % placeholder for frame 11: kernel2
    \ifnum\iFrame=11
        \node[draw=red, very thick, dashed, inner sep=5pt, fit=(kernel2)] {};
    \fi
    
    % placeholder for frame 12: output
    \ifnum\iFrame=12
        \node[draw=red, very thick, dashed, inner sep=5pt, fit=(output)] {};
    \fi

\end{tikzpicture}
}
\end{animateinline}
\end{document}